\chapter{Apprendre et expliquer : le laboratoire d'intelligence artificielle}
\label{ch:apprentissage}

\questionchapitre{Quel dispositif faut-il pour que des expériences d'apprentissage profond sur
des séries hydrologiques soient comparables et reproductibles, et pour qu'on puisse dire ensuite
sur quoi leurs prévisions reposent ?}

\section{Le problème posé}

Prévoir un niveau de nappe est un problème de série temporelle multivariée avec covariables
exogènes : la cible est le niveau, les entrées sont la précipitation, la température et
l'évapotranspiration. Les difficultés propres au domaine sont connues : inertie très variable
d'un aquifère à l'autre, non-stationnarité liée aux prélèvements, lacunes d'observation,
événements extrêmes rares donc sous-représentés dans l'apprentissage.

La littérature propose de nombreuses architectures ; aucune ne s'impose a priori. La question
pratique n'est donc pas « quel modèle choisir » mais « comment mettre en place un dispositif
qui permette de trancher, et de rejouer la comparaison quand les données changent ».

\section{Un catalogue d'architectures sous interface unique}

La plateforme donne accès à une douzaine d'architectures de prévision profonde issues de la
bibliothèque Darts~\cite{darts}, couvrant les grandes familles de l'état de l'art :

\begin{itemize}
  \item \textbf{à mécanisme d'attention} : \emph{temporal fusion transformer}~\cite{lim2021tft},
  transformeur classique ;
  \item \textbf{à décomposition par blocs} : N-BEATS, N-HiTS ;
  \item \textbf{récurrentes} : LSTM, GRU, réseaux récurrents par blocs ;
  \item \textbf{convolutives} : réseaux convolutifs temporels ;
  \item \textbf{récentes à faible coût} : TiDE, TSMixer, DLinear, NLinear.
\end{itemize}

L'inclusion des dernières est délibérée : les modèles linéaires ou quasi linéaires constituent
des références exigeantes sur les séries temporelles, et une architecture lourde qui ne les
dépasse pas doit être écartée. Disposer de la référence faible dans le même outil que le modèle
sophistiqué évite la comparaison biaisée par des conditions expérimentales différentes.

Une fabrique de modèles instancie dynamiquement l'architecture demandée en validant ses
hyperparamètres, ce qui permet d'ajouter une architecture sans modifier l'interface ni l'API.
La figure~\ref{fig:ia-config} montre l'écran où une expérience se déclare.

\begin{figure}[htbp]
\centering
\vue[0.44]{ia-configuration}
\caption{La déclaration d'un entraînement. L'écran expose ce qui doit rester explicite pour
qu'une expérience soit rejouable : architecture, jeu de données horodaté avec sa période et son
nombre de lignes, covariables retenues, fonction de perte, fenêtre d'entrée et horizon de sortie,
partage entre entraînement et validation. Le préréglage métier, en tête, permet de ne rien régler
du tout.}
\label{fig:ia-config}
\end{figure}

\section{Le protocole expérimental}

Le dispositif vise trois propriétés.

La traçabilité, d'abord. Chaque exécution d'entraînement est enregistrée dans MLflow~\cite{mlflow} : paramètres, métriques, artefacts, modèle produit. Les jeux de
données préparés et les modèles entraînés font l'objet de registres distincts, de sorte qu'un
résultat puisse être rattaché sans ambiguïté au couple exact qui l'a produit. Sans ce mécanisme,
une comparaison entre deux expériences menées à quelques semaines d'intervalle n'est pas
défendable.

La séparation entre l'entraînement et l'interface, ensuite. Le code d'entraînement n'a aucune
connaissance de l'interface qui l'observe : il écrit son avancement dans un fichier d'état, que
la couche de présentation lit indépendamment et retransmet au navigateur par un flux
d'événements. Cette séparation, visible sur la figure~\ref{fig:archi-plateforme}, a une
conséquence pratique importante : un entraînement lancé depuis un script se suit dans
l'interface, et réciproquement.

La conduite d'expériences longues, enfin. L'entraînement s'appuie sur PyTorch Lightning~\cite{pytorch_lightning} et s'effectue sur
processeur graphique lorsque celui-ci est disponible, avec arrêt anticipé et suivi de
l'évolution des pertes en apprentissage et en validation. L'utilisateur voit progresser son
entraînement au lieu d'attendre un résultat opaque, et peut l'interrompre s'il diverge.

Le modèle produit conserve ensuite ce qui l'a fait (figure~\ref{fig:ia-modele}), de sorte qu'une
prévision se rattache sans ambiguïté au couple exact qui l'a produite.

\begin{figure}[htbp]
\centering
\vue{ia-modele}
\caption{Un modèle entraîné, tel que la page de prévision le présente : architecture, station,
jeu de données horodaté, variables employées, métriques de test, et la fenêtre d'entrée comme
l'horizon avec lesquels il a été appris. Un modèle dont on ne saurait pas dire ces choses ne
serait pas comparable à un autre. \lecas}
\label{fig:ia-modele}
\end{figure}

\section{Optimisation d'hyperparamètres}

La comparaison d'architectures n'a de sens que si chacune est évaluée dans une configuration
raisonnable ; comparer un modèle réglé à un modèle laissé aux valeurs par défaut ne renseigne
que sur le réglage. Optuna~\cite{optuna} est donc intégré pour la recherche automatisée d'hyperparamètres, avec un
échantillonnage guidé par les essais déjà évalués. Chaque essai est par ailleurs interrompu dès
que sa perte de validation cesse de progresser. L'élagage entre essais, qui abandonnerait une
configuration en cours au vu des premières époques, n'est pas activé : ce serait le premier
réglage à ajouter pour réduire le coût d'une campagne.

\section{Rendre l'outil utilisable par des non-spécialistes}

Un dispositif d'expérimentation complet est, par nature, exigeant : longueur de fenêtre
d'entrée, taille de lot, nombre d'époques, patience d'arrêt anticipé. Ces réglages sont un
obstacle réel pour un hydrogéologue qui souhaite simplement obtenir une prévision.

Des \textbf{configurations prêtes à l'emploi} ont donc été définies par cas d'usage métier
plutôt que par architecture : prévision piézométrique à court terme (sept jours), à moyen terme
(trente jours), à long terme (quatre-vingt-dix jours), et prévision de débit à quatorze jours.
Chacune fixe une architecture et un jeu d'hyperparamètres cohérents, choisis à partir des
recommandations de la bibliothèque et des références publiées sur séries piézométriques.

Ces configurations sont des points de départ raisonnables, non des réglages optimisés pour une
station donnée. La documentation le dit explicitement et oriente vers la recherche automatisée
d'hyperparamètres pour aller plus loin. Il aurait été facile, et trompeur, de les présenter
comme optimales.

Les échanges avec les hydrogéologues ont d'ailleurs montré que cette réserve n'est pas de pure
forme. La réponse d'un aquifère à la pluie varie fortement selon sa nature : certains systèmes
sont inertiels et réagissent avec une forte latence, d'autres se rechargent presque
directement. Les cycles observés vont de l'annuel au pluriannuel. La profondeur d'historique à
fournir en entrée, l'horizon de prévision pertinent et le pas d'agrégation devraient donc être
choisis selon le type de nappe, comme le sont déjà les configurations de modèles métier du
chapitre~\ref{ch:metier}. Cette adaptation par famille d'aquifère n'a pas été faite pour les
modèles appris : elle constitue une piste identifiée, et vraisemblablement le premier gain à
aller chercher.

\section{Lire une prévision avant de l'expliquer}

Le dispositif expose la prévision de deux manières complémentaires, dont la seconde conditionne
ce qu'on peut dire de la première.

La première est la \textbf{fenêtre glissante} sur le jeu de test
(figure~\ref{fig:ia-jeu-de-test}). La fenêtre d'entrée et l'horizon se déplacent ensemble sur
toute la période réservée, ce qui évite de conclure d'une fenêtre unique et bien choisie.

\begin{figure}[htbp]
\centering
\vue{ia-jeu-de-test}
\caption{Le jeu de test entier, que le modèle n'a pas vu à l'entraînement. La zone claire est la
fenêtre d'entrée, la bande ocre l'horizon de prévision, et les deux se positionnent n'importe où
sur la période. La vue d'ensemble sert d'abord à situer la fenêtre examinée dans la dynamique de
l'ouvrage, saison humide ou étiage.}
\label{fig:ia-jeu-de-test}
\end{figure}

La seconde est la \textbf{fenêtre examinée de près}, avec ses écarts chiffrés
(figure~\ref{fig:prevision}). L'erreur y est rapportée à l'amplitude de l'ouvrage, et non
seulement exprimée en mètres : une erreur de quelques centimètres n'a pas le même sens sur une
nappe dont le battement annuel est décimétrique et sur une nappe dont il dépasse le mètre. C'est
cette mise en rapport, plus que l'erreur brute, qui permet à un hydrogéologue de dire si la
prévision lui sert.

\begin{figure}[htbp]
\centering
\vue[0.62]{prevision-fenetre}\hfill\vue[0.33]{prevision-metriques}
\caption{Une fenêtre de prévision confrontée aux observations, et les écarts associés. La
dernière ligne du panneau rapporte l'erreur absolue moyenne et l'erreur quadratique à l'écart
interquartile de la série, ce qui les rend comparables d'un ouvrage à l'autre. La précision
directionnelle indique par ailleurs la part des pas où le sens de variation est correct, qui est
souvent ce qui intéresse un gestionnaire. \lecas}
\label{fig:prevision}
\end{figure}

\section{Analyser les prévisions produites}
\label{sec:xai}

Une prévision de niveau de nappe n'est mobilisable pour décider que si l'on peut dire sur quoi
elle repose. Les décisions publiques rappelées au chapitre~\ref{ch:contexte} doivent pouvoir être
motivées, et l'interlocuteur d'un modèle de nappe est un hydrogéologue qui n'accordera sa
confiance qu'à des explications compatibles avec sa compréhension du système. Le laboratoire
outille donc trois questions : \emph{sur quoi le modèle s'appuie-t-il}, \emph{que se
passerait-il dans d'autres conditions}, et \emph{cette chronique est-elle seulement explicable
par le climat}.

La première question mobilise cinq familles d'attribution complémentaires :
l'importance par perturbation, les valeurs de Shapley~\cite{lundberg2017shap} et leur variante
pour séries temporelles~\cite{bento2021timeshap}, les attributions par
gradients~\cite{sundararajan2017ig,captum}, les poids d'attention pour les architectures qui en
possèdent, et la décomposition du signal en tendance, saisonnalité et résidu. Les réunir dans un
même outil ne vise pas l'exhaustivité mais la \textbf{confrontation} : des méthodes reposant sur
des principes différents et qui convergent donnent une explication crédible, tandis que leur
divergence signale une explication fragile, information qu'une méthode unique aurait masquée.

Quatre de ces cinq familles sont atteignables depuis l'interface, dans le panneau
d'explicabilité de la page de prévision : l'importance par perturbation, les valeurs de Shapley,
les attributions par gradients et la décomposition du signal. Les poids d'attention font
exception. Le calcul existe dans la bibliothèque, le point d'accès existe dans l'API et le
client web sait l'appeler, mais aucun composant de l'interface ne le sollicite : ils ne sont donc
pas accessibles à un utilisateur.

La deuxième question relève des explications contrefactuelles, formalisées par Wachter
\emph{et al.}~\cite{wachter2018counterfactual} et transposées depuis peu à la prévision de
séries temporelles~\cite{wang2023forecastcf}. Elle a demandé un développement propre, parce que
la réponse directe ne vaut rien ici. Laisser un algorithme modifier librement les entrées jusqu'à ce que la prévision
atteigne la cible produit des scénarios corrects au sens de l'optimisation et physiquement
absurdes : de l'air réchauffé sans évaporation supplémentaire, une pluie incohérente d'un jour
au suivant, une saison qui n'existe nulle part. Les hydrogéologues y ont substitué deux
exigences, le respect du bilan de masse et la cohérence avec les échelles de temps auxquelles un
aquifère réagit.

\subsection{Contrefactuels sous contrainte physique}
\label{sec:physcf}

L'implémentation retenue, désignée \physcf{}, a été conçue au cours de ces travaux ; un
prototype séparé, décrit en section~\ref{sec:travaux-perso}, sert à l'éprouver hors de
l'application. Elle traduit ces exigences en renonçant à perturber les entrées valeur par
valeur, et ne règle que sept paramètres décrivant un régime climatique : un facteur de pluie par
saison, un décalage de température, un décalage du calendrier, et un résidu d'évaporation
étroitement borné. L'évapotranspiration n'est pas libre, elle est calculée à partir de la
température au taux d'environ \SI{7}{\percent} par degré, ordre de grandeur tiré de la relation
de Clausius-Clapeyron. Le scénario reste donc cohérent par construction plutôt que redressé
après coup, et la réponse s'énonce dans les termes du métier : \emph{il aurait fallu des hivers
un cinquième plus humides et un climat plus frais d'un degré}. Une seconde méthode, adaptée de
CoMTE~\cite{ates2021comte}, répond à la même question en allant chercher dans l'historique une
période réellement observée plutôt qu'en construisant un climat, ce qui donne un point de
comparaison de nature différente. Chaque scénario produit est enfin rejoué dans un modèle métier
calibré indépendamment, celui du chapitre~\ref{ch:metier} : s'il fait diverger les deux modèles
bien plus qu'ils ne divergeaient sur la situation observée, il est rejeté.

Ces deux méthodes et leur validation croisée sont livrées dans la bibliothèque métier et
exposées par l'API, où elles disposent chacune de leur point d'accès. Le prototype de recherche
dont \physcf{} est issue vit dans un dépôt distinct, décrit plus loin. Elles ne sont en revanche
\textbf{pas atteignables depuis l'interface}. La page qui les pilote existe dans le dépôt, avec
son sélecteur de cible et sa vue de comparaison, mais elle n'est déclarée dans aucune route de
l'application : le laboratoire d'intelligence artificielle propose quatre onglets, données,
entraînement, prévision et modèles, et aucun d'eux ne mène aux contrefactuels. Un utilisateur ne
peut donc en produire aucun sans passer par l'API ou par un carnet de calcul. Rebrancher cette
page est un travail court, puisque tout ce qu'elle appelle fonctionne, mais il n'a pas été fait
et il conditionne tout usage par un non-développeur.

La troisième question renverse la perspective. Une chronique qu'aucun modèle climatique ne
reproduit n'est pas nécessairement mal modélisée : elle peut porter la trace d'un prélèvement,
que le forçage climatique ne contient pas. Le problème est non supervisé, faute d'historique
étiqueté. La démarche part des résidus d'un modèle métier calibré pour isoler des périodes
tenues pour non perturbées, puis croise deux lectures indépendantes dont les verdicts sont
fusionnés. L'explicabilité y joue un rôle inhabituel : elle ne justifie pas une prédiction, elle
détecte une dérive de ce sur quoi le modèle s'appuie, signe que le régime du système a changé.
L'annotation qui en résulte sert à écarter les stations impropres à l'apprentissage, et figure
parmi les cas d'usage identifiés avec le BRGM.

\section{Deux travaux personnels conduits en marge}
\label{sec:travaux-perso}

Deux questions soulevées par la plateforme ne pouvaient pas être traitées avec elle. Toutes deux
demandaient de sortir de l'application pour isoler la question de l'outil qui la pose, et toutes
deux ont donné lieu à un développement séparé, mené à titre personnel dans le cadre de ces
travaux. Elles figurent ici parce qu'elles éclairent des choix décrits plus haut, et parce qu'un
repreneur qui ignorerait la seconde referait une erreur qu'elle documente.

\subsection{Éprouver la méthode contrefactuelle hors de l'application}

La méthode décrite en section~\ref{sec:physcf} devait être évaluée, et la plateforme ne s'y
prête pas : elle mêle l'instrument et le cas, alors qu'une évaluation demande un petit nombre de
sites choisis, des modèles entraînés pour eux et des mesures répétables. Le prototype
\physcf{} a donc été construit à part, autour de quatre stations couvrant le gradient qui va de
l'aquifère le plus réactif au plus inertiel : un alluvial de Moselle, un karst du Val d'Orléans,
une craie de la Somme et un calcaire de Beauce.

Le but est de vérifier une hypothèse précise : que le choix de \emph{l'espace} de perturbation
compte davantage que celui de l'algorithme d'optimisation. Une méthode de référence opérant en
espace brut, sur un millier de dimensions avec des pénalités molles, est comparée à \physcf{}, qui
optimise quatre à sept paramètres dans un espace où la cohérence physique est vraie par
construction. Chaque scénario est confronté à un modèle métier calibré indépendamment.

Ce qui est acquis est le dispositif complet : les opérateurs de perturbation, l'optimisation par
gradient et par recherche bayésienne, la validation croisée par le modèle métier, les cibles
définies sur les classes de l'indicateur piézométrique, et la suite de tests qui les couvre. Le
protocole d'évaluation est écrit lui aussi, jusqu'aux mesures à relever et aux tableaux qui
doivent les recevoir. Ces tableaux sont vides : la campagne n'a pas été conduite, et le
chapitre~\ref{ch:bilan} dit ce qu'elle demande d'abord.

\subsection{Un banc d'essai des représentations de chroniques piézométriques}

La détection des chroniques perturbées par un prélèvement, évoquée en section~\ref{sec:xai},
prévoit un étage qui suit la dérive d'une représentation apprise de la chronique. Rien ne disait
quel encodeur employer, et la littérature ne comparait aucune de ces méthodes sur des données
hydrogéologiques. La question s'est révélée plus large que son motif : elle porte sur ce qu'une
compression apprise retient d'une chronique piézométrique.

Sept encodeurs ont été comparés sur \num{5116} stations, avec un partage de validation croisée
qui tient chaque département entier du même côté : aucune station n'est ainsi évaluée par un
modèle qui aurait appris sur ses voisines immédiates, dont la chronique lui ressemble par le
seul effet de la proximité. Le protocole principal est \textbf{univarié} : chaque encodeur ne reçoit que le
niveau. Une variante y ajoute les trois covariables climatiques, non comme un raffinement mais
comme une expérience de contrôle.

\begin{table}[H]
\centering
\caption{Classification du milieu hydrogéologique à partir du seul niveau, précision équilibrée
sous partage spatial. La dernière colonne mesure ce que la représentation retient de la
latitude, quantité qui n'a aucune raison d'aider à reconnaître un aquifère.}
\label{tab:banc-plongements}
\begin{tabularx}{\textwidth}{@{}Yrrl@{}}
\toprule
\textbf{Encodeur} & \textbf{Précision équilibrée} & \textbf{$R^2$ latitude} & \textbf{Nature} \\
\midrule
Signatures hydrologiques & \num{0,303} & \num{-0,021} & 29 descripteurs métier \\
Catch22                  & \num{0,259} & \num{0,117}  & statistiques génériques \\
MultiRocket              & \num{0,256} & \num{0,365}  & convolutions aléatoires \\
SoftCLT                  & \num{0,241} & \num{0,222}  & contrastif appris \\
TS2Vec                   & \num{0,233} & \num{0,174}  & contrastif appris \\
Analyse en composantes   & \num{0,170} & \num{0,003}  & référence linéaire \\
Bruit gaussien           & \num{0,148} & \num{-0,012} & référence basse \\
\bottomrule
\end{tabularx}
\end{table}

Trois résultats en sortent. Le premier est que \textbf{vingt-neuf descripteurs écrits à la main
devancent l'apprentissage contrastif} sur le signal endogène, avec un écart statistiquement
significatif contre les deux encodeurs contrastifs et un écart non significatif contre les deux
méthodes génériques restantes, Catch22 et MultiRocket. Une recherche d'hyperparamètres dédiée à
l'un des encodeurs contrastifs n'en comble que \numrange{0,015}{0,020}, le reste étant structurel
plutôt qu'un défaut de réglage. Les mêmes descripteurs prédisent les signatures physiques d'une
chronique avec un coefficient de détermination proche de \num{0,89}, contre
\numrange{0,08}{0,19} pour les représentations apprises. Cet avantage tient cependant à la tâche
et ne se transporte pas : sur une régression de tendance longue, ce sont les descripteurs
génériques qui l'emportent, tandis que les descripteurs métier retombent au niveau d'une
prédiction faite à partir de la seule position géographique. Aucun encodeur ne domine partout,
et le choix suit l'usage qu'on veut en faire.

Le deuxième porte sur la validation. Une validation croisée qui ignore la structure spatiale
porte la valeur de \num{0,303} à \num{0,376} pour le même encodeur : l'autocorrélation
géographique des étiquettes suffit à gonfler une performance sans qu'aucune erreur de calcul ne
soit commise.

Le troisième est le plus utile ici, et le plus inattendu. \textbf{Ajouter les covariables
climatiques brutes améliore la classification, et pour la mauvaise raison.} La précision
équilibrée gagne de \numrange{0,059}{0,146} selon l'encodeur, mais la latitude devient au même
moment prédictible à partir de la représentation, avec un gain de \numrange{0,488}{0,633} sur le
même coefficient. Deux stations voisines partagent leur maille climatique ; le forçage écrit donc
dans la représentation une structure géographiquement lisse, et l'essentiel du gain apparent est
un transfert de cette structure, non un progrès de discrimination hydrogéologique.

Ce résultat conditionne l'interprétation des modèles multivariés de ce chapitre, qui consomment
exactement ces trois covariables. Une réserve accompagne cependant sa transposition : le banc
travaille sur une grille au quart de degré, alors que l'entrepôt en produit une au dixième de
degré, dont la structure est moins lisse. Le mécanisme vaut, son ampleur reste à mesurer sur la
chaîne décrite au chapitre~\ref{ch:entrepot}. Le remède proposé, employer des anomalies ou une
recharge dérivée du modèle métier plutôt que le forçage brut, figure parmi les chantiers du
chapitre~\ref{ch:bilan}.

Le banc explique enfin pourquoi l'étage de dérive de représentation reste inerte dans la
plateforme : les encodeurs qu'il attend ne se sont pas montrés supérieurs, sur le signal
endogène, à des descripteurs que la bibliothèque métier calcule déjà.

Une précision s'impose sur l'état du dépôt, puisqu'il est public et que le tableau ci-dessus
invite à s'y reporter. Les campagnes multivariées, la régression de tendance et les recherches
d'hyperparamètres y sont versionnées avec leurs métriques ; la campagne univariée dont ce
tableau est tiré ne l'est pas, et les carnets de vérification n'ont pas été exécutés. Les
valeurs rapportées ici ont été recoupées avec les artefacts disponibles, mais les rendre
rejouables d'un bout à l'autre reste à faire, et figure au chapitre~\ref{ch:bilan}.

\section{Ce qui est établi, et ce qui ne l'est pas}

\textbf{Ce qui est acquis} est le dispositif : un environnement d'expérimentation instrumenté,
traçable, qui rend les comparaisons possibles et rejouables, accessible à des utilisateurs de
niveaux d'expertise différents, et dont le cœur est réutilisable hors de l'application.

\textbf{Ce qui n'est pas établi dans ce rapport} est le classement des architectures \emph{de
prévision} sur le cas hydrogéologique. La section précédente compare des encodeurs de
représentation, question distincte traitée par un dispositif distinct ; sur la prévision
elle-même, ce rapport ne produit aucun classement. Les campagnes d'entraînement conduites au
cours des travaux ont servi à valider le dispositif et à orienter les développements ; les jeux
de données préparés et les points de sauvegarde ont été purgés des dépôts lors de la passation.
Reconstituer a posteriori un tableau de performances de prévision à partir de souvenirs ne
serait pas un résultat scientifique. La
campagne d'évaluation comparative, qui opposerait les modèles appris aux modèles métier du
chapitre~\ref{ch:metier} sur un échantillon représentatif de stations et de familles
d'aquifères, reste donc à conduire, et l'outillage nécessaire est en place pour la mener.

Cela vaut au même titre pour les outils d'analyse. Ils mettent en œuvre l'exigence formulée par
les hydrogéologues, et le dispositif prévoit de quoi les comparer entre eux comme de quoi
vérifier leurs résultats sur un modèle indépendant. Ils sont opérationnels au niveau de la
bibliothèque et de l'API ; pour les contrefactuels, ils ne le sont pas au niveau de l'interface,
faute du branchement signalé plus haut. Mais aucune évaluation
quantitative n'a été conduite, sur aucun échantillon, selon aucune métrique de qualité
d'explication établie. Ces développements sont \textbf{préliminaires}, et c'est leur objet :
la question de ce qu'est une bonne explication relève des travaux en cours évoqués au
chapitre~\ref{ch:collaborations}, auxquels ce laboratoire fournit son terrain.

Une limite d'ingénierie reste à corriger. La suite de tests de la plateforme compte
\num{552} tests, dont \num{547} passent (mesuré le 24 août 2026). Aucun ne s'exécute
automatiquement : les chaînes d'intégration des deux dépôts ne comportent que des étapes de
construction et de déploiement, sans étape de test. Les cinq cas non passants, quatre
échecs et un test ignoré, sont documentés et aucun ne signale un défaut applicatif, mais la
purge des journaux se trouve de fait non testée. C'est le premier correctif à apporter par toute
personne reprenant ces travaux.
